\documentclass[12pt]{article}
\usepackage[utf8]{inputenc}
\usepackage[T1]{fontenc}
\usepackage[margin=1in]{geometry}
\usepackage{setspace}
\usepackage{xcolor}
\usepackage[colorlinks=true, urlcolor=blue!60!black, linkcolor=blue!60!black]{hyperref}
\setlength{\parindent}{0em}
\setlength{\parskip}{0.6em}
\renewcommand{\baselinestretch}{1.0}
\begin{document}
\begin{center}
{\LARGE \textit{Gibbons v. Ogden}, 22 U.S. (9 Wheat.) 1 (1824)\par}
\vspace{0.5em}
{\large Case Brief --- Two-Page Edition\par}
\end{center}

\section{Memory Jogger}
Congress's commerce power covers interstate navigation and is plenary, so a federal coasting license preempts New York's steamboat monopoly (\href{https://www.courtlistener.com/opinion/85412/gibbons-v-ogden/}{22 U.S.\ 1}).

\section{Detailed Case Facts}
New York granted Livingston and Fulton an exclusive right to run steamboats in its waters; Aaron Ogden held the franchise for the Elizabethtown--New York City route. Thomas Gibbons ran rival steamboats licensed under the federal Coasting Act of 1793. Neighboring states had retaliated against the monopoly. Ogden sued in New York chancery to enjoin Gibbons.

\section{Procedural History}
Chancellor Kent enjoined Gibbons, and New York's highest court affirmed. After a first appeal was dismissed for want of a final decree, the Supreme Court took the case.

\section{Judicial Votes}
Reversed, 6--0. Marshall wrote for Washington, Todd, Duvall, and Story; Johnson concurred in the judgment; Thompson did not participate.

\section{Holding}
Commerce includes navigation, and commerce ``among the several States'' reaches into a state's interior, though not commerce ``completely internal'' to one state. Congress's power is plenary, and the federal license preempts the conflicting monopoly under the Supremacy Clause.

\section{Analysis of Opinions}
Marshall rejected strict construction and defined commerce as ``commercial intercourse'' in all its branches; ``among'' means ``intermingled with''; and ``to regulate'' means to prescribe the rule, a power checked mainly by elections. State inspection and health laws rest on the police power. Declining to decide Webster's exclusivity theory, Marshall held that the Coasting Act license authorized Gibbons's trade, voiding the monopoly to the extent of the conflict. Johnson would have held the commerce power exclusive, voiding the monopoly even without federal legislation.

\section{Examples: Future Applications}
\textbf{Same-Side:} a state's exclusive river-ferry franchise yields to a federal interstate license. \textbf{Opposite-Side:} a local livestock-inspection law with no conflicting federal statute is valid police-power legislation. \textbf{Fence-Sitter:} an intrastate Hudson ferry monopoly carrying many out-of-state travelers.

\section{Critique}
Barnett finds Marshall's ``commerce'' faithful to original meaning but not to later expansions; Balkin reads it more broadly. Norman Williams faults Marshall's statutory ``hedge'' for leaving the dormant commerce question confused until \textit{Cooley}.

\section{Key Quotations}
\begin{quote}
``It is the power to regulate; that is, to prescribe the rule by which commerce is to be governed. This power, like all others vested in Congress, is complete in itself, may be exercised to its utmost extent, and acknowledges no limitations, other than are prescribed in the Constitution.''
\end{quote}
\noindent --- Chief Justice Marshall, 22 U.S.\ at 196.

\end{document}
